% Chapter 2.9: Memory Circuits (1-OS) - LaTeX Beamer Slides
% Structured Computer Architecture: An Order-Based Approach
% Authors: John Glossner and Gheorghe M. Ştefan

\documentclass[aspectratio=169,11pt]{beamer}

% Theme and color scheme (simple, print-friendly)
\usetheme{default}
\usecolortheme{default}

% Simple color scheme for printing
\definecolor{titlecolor}{RGB}{0,0,0}
\definecolor{accentcolor}{RGB}{0,0,139}
\definecolor{textcolor}{RGB}{0,0,0}

% Set colors
\setbeamercolor{title}{fg=titlecolor}
\setbeamercolor{frametitle}{fg=titlecolor}
\setbeamercolor{structure}{fg=accentcolor}
\setbeamercolor{normal text}{fg=textcolor}

% Remove navigation symbols
\setbeamertemplate{navigation symbols}{}

% Simple footline
\setbeamertemplate{footline}[frame number]

% Packages
\usepackage{amsmath}
\usepackage{amsfonts}
\usepackage{amssymb}
\usepackage{graphicx}
\usepackage{tikz}
\usepackage{booktabs}
\usepackage{array}
\usepackage{listings}
\usepackage{xcolor}

% Set graphics path for images (relative to slides directory)
\graphicspath{{../Images/}}

% Title information
\title[Chapter 2.9: Memory Circuits]{Structured Computer Architecture\\Chapter 2.9: Memory Circuits\\First-Order Systems (1-OS)}
\subtitle{An Order-Based Approach}
\author{John Glossner and Gheorghe M. Ştefan}
\institute{Based on the textbook}
\date{\today}

% Custom commands
\newcommand{\concept}[1]{\textbf{\color{accentcolor}#1}}
\newcommand{\defterm}[1]{\textit{#1}}

\begin{document}

% Title slide
\begin{frame}
\titlepage
\end{frame}

% Table of contents
\begin{frame}{Outline}
\tableofcontents
\end{frame}

% Learning Objectives
\begin{frame}{Learning Objectives}
By the end of this chapter, you will be able to:
\begin{itemize}
\item Understand the concept of first-order systems (1-OS) and feedback loops
\item Analyze the behavior of latches and flip-flops
\item Design clocked memory circuits and understand timing constraints
\item Explain the operation of RAM and register files
\item Distinguish between synchronous and asynchronous memory systems
\item Understand the trade-offs in memory circuit design
\end{itemize}
\end{frame}

% Section 1: Introduction to 1-OS
\section{Introduction to First-Order Systems}

\begin{frame}{From 0-OS to 1-OS: Adding Memory}
\begin{itemize}
\item \concept{0-OS Limitation}: Output depends only on current inputs
\begin{itemize}
\item No memory of previous states
\item Cannot store information
\item Purely reactive behavior
\end{itemize}
\vspace{0.5cm}
\item \concept{1-OS Innovation}: Introduction of the first feedback loop
\begin{itemize}
\item Output fed back to input
\item Creates memory capability
\item Enables state storage
\item Partial autonomy achieved
\end{itemize}
\vspace{0.3cm}
\item \concept{Key Principle}: 
\begin{itemize}
\item Memory function obtained through reverse connection
\item Circuit maintains state according to command
\item Foundation for all sequential logic
\end{itemize}
\end{itemize}
\end{frame}

\begin{frame}{The Feedback Loop Concept}
\begin{itemize}
\item \concept{Feedback Mechanism}:
\begin{itemize}
\item Output signal returns to influence input
\item Creates dependency on previous state
\item Enables information storage
\end{itemize}
\vspace{0.5cm}
\item \concept{State Autonomy}:
\begin{itemize}
\item State can be maintained without continuous input
\item Temporary input can have permanent effect
\item System "remembers" past events
\end{itemize}
\vspace{0.5cm}
\item \concept{Applications in Computing}:
\begin{itemize}
\item Register storage in processors
\item Memory cells in RAM
\item State machines in control units
\item Pipeline registers
\end{itemize}
\end{itemize}
\end{frame}

% Section 2: Elementary Latches
\section{Elementary Latches}
\begin{frame}{The Simplest Memory Element}
\begin{columns}
\begin{column}{0.6\textwidth}
\concept{Latch Function}: Store temporary signal transitions
\begin{itemize}
\item Capture transition from 1 to 0 (reset-only latch)
\item Capture transition from 0 to 1 (set-only latch)
\item Maintain state indefinitely
\end{itemize}
\vspace{0.3cm}
\concept{Basic Principle}:
\begin{itemize}
\item Feedback loop from output to input
\item Signal propagation through circuit
\item State becomes self-sustaining
\end{itemize}
\vspace{0.3cm}
\concept{Limitation}: 
\begin{itemize}
\item Simple latches cannot "forget"
\item Need combination for full functionality
\item Asymmetric command signals
\end{itemize}
\end{column}
\begin{column}{0.4\textwidth}
\begin{center}
\textbf{Elementary Latches}
\vspace{0.3cm}

\tikz{
\node[draw, rectangle, minimum width=2cm, minimum height=1cm] (and) at (0,1) {AND Loop};
\node[draw, rectangle, minimum width=2cm, minimum height=1cm] (or) at (0,0) {OR Loop};
\node[draw, rectangle, minimum width=2cm, minimum height=1cm] (sr) at (0,-1) {Set-Reset};
\draw[->] (and) -- (or);
\draw[->] (or) -- (sr);
}
\end{center}
\end{column}
\end{columns}
\end{frame}

\begin{frame}{AND and OR Loop Latches}
\begin{columns}
\begin{column}{0.5\textwidth}
\concept{AND Loop (Reset-Only)}:
\begin{itemize}
\item Output connected to one AND input
\item Reset' pulse applied to other input
\item Captures transition to 0
\item Output fixed at 0 after reset
\end{itemize}

\vspace{0.5cm}
\concept{Behavior}:
\begin{itemize}
\item Normal state: Output = 1
\item Reset' = 0: Forces output to 0
\item Output remains 0 permanently
\item Cannot return to 1
\end{itemize}
\end{column}
\begin{column}{0.5\textwidth}
\concept{OR Loop (Set-Only)}:
\begin{itemize}
\item Output connected to one OR input
\item Set pulse applied to other input
\item Captures transition to 1
\item Output fixed at 1 after set
\end{itemize}

\vspace{0.5cm}
\concept{Behavior}:
\begin{itemize}
\item Normal state: Output = 0
\item Set = 1: Forces output to 1
\item Output remains 1 permanently
\item Cannot return to 0
\end{itemize}
\end{column}
\end{columns}
\end{frame}

\begin{frame}{Set-Reset (SR) Latch}
\begin{itemize}
\item \concept{Combination Solution}: Merge set-only and reset-only latches
\begin{itemize}
\item Provides both set and reset capability
\item Can switch between two stable states
\item Fundamental memory structure
\end{itemize}
\vspace{0.5cm}
\item \concept{Operation}:
\begin{itemize}
\item \textbf{Set command (S)}: Active high, sets output to 1
\item \textbf{Reset command (R')}: Active low, resets output to 0
\item \textbf{Asymmetric commands}: Different activation levels
\end{itemize}
\vspace{0.3cm}
\item \concept{State Table}:
\begin{center}
\begin{tabular}{|c|c|c|l|}
\hline
S & R' & Q & Action \\
\hline
0 & 1 & Q & Hold \\
1 & 1 & 1 & Set \\
0 & 0 & 0 & Reset \\
1 & 0 & ? & Forbidden \\
\hline
\end{tabular}
\end{center}
\end{itemize}
\end{frame}

\begin{frame}{Symmetric SR Latches}
\begin{columns}
\begin{column}{0.5\textwidth}
\concept{NAND SR Latch}:
\begin{itemize}
\item Uses De Morgan's law transformation
\item Both commands active low (S', R')
\item More common in practice
\item Symmetric command structure
\end{itemize}

\vspace{0.3cm}
\concept{Truth Table}:
\begin{center}
\begin{tabular}{|c|c|c|}
\hline
S' & R' & Q \\
\hline
1 & 1 & Hold \\
0 & 1 & 1 \\
1 & 0 & 0 \\
0 & 0 & Invalid \\
\hline
\end{tabular}
\end{center}
\end{column}
\begin{column}{0.5\textwidth}
\concept{NOR SR Latch}:
\begin{itemize}
\item Alternative De Morgan transformation
\item Both commands active high (S, R)
\item Symmetric high-active commands
\item Equivalent functionality
\end{itemize}

\vspace{0.3cm}
\concept{Truth Table}:
\begin{center}
\begin{tabular}{|c|c|c|}
\hline
S & R & Q \\
\hline
0 & 0 & Hold \\
1 & 0 & 1 \\
0 & 1 & 0 \\
1 & 1 & Invalid \\
\hline
\end{tabular}
\end{center}
\end{column}
\end{columns}

\vspace{0.3cm}
\begin{center}

\end{center}
\end{frame}

% Section 3: Clocked Latches
\section{Clocked Latches and Timing}

\begin{frame}{The First Latch Problem}
\begin{itemize}
\item \concept{Problem Statement}: 
\begin{itemize}
\item Same inputs specify \textbf{how} and \textbf{when} latch switches
\item Immediate reactivity to input changes
\item Susceptible to noise and glitches
\item Mechanical switch bouncing issues
\end{itemize}
\vspace{0.5cm}
\item \concept{Consequences}:
\begin{itemize}
\item Unpredictable state changes
\item Noise-induced switching
\item Difficult timing control
\item Unreliable operation
\end{itemize}
\vspace{0.5cm}
\item \concept{Solution Approach}:
\begin{itemize}
\item Separate "how" from "when"
\item Use clock signal for timing control
\item Condition inputs with clock
\item Achieve synchronous operation
\end{itemize}
\end{itemize}
\end{frame}
\begin{frame}{Clocked SR Latch}
\begin{columns}
\begin{column}{0.6\textwidth}
\concept{Design Solution}: Add clock conditioning
\begin{itemize}
\item S and R inputs specify \textbf{how} to switch
\item Clock (CK) input specifies \textbf{when} to switch
\item Decoupling of timing from function
\item Synchronous state changes
\end{itemize}
\vspace{0.3cm}
\concept{Implementation}:
\begin{itemize}
\item Two additional NAND gates
\item S and R conditioned by clock
\item Latch responds only when CK = 1
\item Immune to input changes when CK = 0
\end{itemize}
\vspace{0.3cm}
\concept{Operation}:
\begin{itemize}
\item \textbf{CK = 0}: Latch holds current state
\item \textbf{CK = 1}: Latch responds to S and R
\item Eliminates noise sensitivity
\item Provides controlled timing
\end{itemize}
\end{column}
\begin{column}{0.4\textwidth}
\begin{center}

\end{center}
\end{column}
\end{columns}
\end{frame}}
\begin{frame}{D Latch (Data Latch)}
\begin{columns}
\begin{column}{0.6\textwidth}
\concept{Motivation}: Eliminate forbidden state of SR latch
\begin{itemize}
\item Single data input instead of S and R
\item No invalid input combinations
\item Simpler interface
\end{itemize}
\vspace{0.3cm}
\concept{Implementation}:
\begin{itemize}
\item D input connected to S
\item $\overline{D}$ connected to R through inverter
\item Clock enables data transfer
\item Transparent when clock is high
\end{itemize}
\vspace{0.3cm}
\concept{Behavior}:
\begin{itemize}
\item \textbf{CK = 1}: Q follows D (transparent)
\item \textbf{CK = 0}: Q holds last value
\item No forbidden states
\item Simple data storage
\end{itemize}
\end{column}
\begin{column}{0.4\textwidth}
\begin{center}

\end{center}
\end{column}
\end{columns}
\end{frame}

% Section 4: Flip-Flops
\section{Flip-Flops and Edge-Triggered Circuits}

\begin{frame}{From Level to Edge Triggering}
\begin{itemize}
\item \concept{Latch Limitation}: Level-sensitive operation
\begin{itemize}
\item Transparent while clock is active
\item Multiple changes during clock high
\item Timing window issues
\end{itemize}
\vspace{0.5cm}
\item \concept{Flip-Flop Advantage}: Edge-triggered operation
\begin{itemize}
\item Changes only on clock edge
\item Single transition per clock cycle
\item Precise timing control
\item Better noise immunity
\end{itemize}
\vspace{0.3cm}
\item \concept{Implementation Methods}:
\begin{itemize}
\item Master-slave configuration
\item Edge detection circuits
\item Pulse-triggered designs
\end{itemize}
\end{itemize}
\end{frame}
\begin{frame}{Master-Slave D Flip-Flop}
\begin{columns}
\begin{column}{0.6\textwidth}
\concept{Structure}: Two latches in series
\begin{itemize}
\item Master latch: Enabled when CLK = 1
\item Slave latch: Enabled when CLK = 0
\item Inverted clock between stages
\end{itemize}
\vspace{0.3cm}
\concept{Operation Sequence}:
\begin{enumerate}
\item \textbf{CLK = 1}: Master captures input D
\item \textbf{CLK = 0}: Slave captures master output
\item Output changes only on falling edge
\item Input isolated from output
\end{enumerate}
\vspace{0.3cm}
\concept{Advantages}:
\begin{itemize}
\item Edge-triggered behavior
\item No transparency issues
\item Predictable timing
\item Standard building block
\end{itemize}
\end{column}
\begin{column}{0.4\textwidth}
\begin{center}

\end{center}
\vspace{0.3cm}
\begin{center}

\end{center}
\end{column}
\end{columns}
\end{frame}}

\begin{frame}{Flip-Flop Types and Applications}
\begin{itemize}
\item \concept{D Flip-Flop}:
\begin{itemize}
\item Data storage and transfer
\item Pipeline registers
\item State machine implementation
\item Most common type
\end{itemize}
\vspace{0.3cm}
\item \concept{JK Flip-Flop}:
\begin{itemize}
\item No forbidden states
\item Toggle capability
\item Counter applications
\item More complex control
\end{itemize}
\vspace{0.3cm}
\item \concept{T Flip-Flop}:
\begin{itemize}
\item Toggle operation
\item Frequency division
\item Counter circuits
\item Simple interface
\end{itemize}
\end{itemize}
\end{frame}

% Section 5: Timing Analysis
\section{Timing Analysis and Constraints}

\begin{frame}{Critical Timing Parameters}
\begin{itemize}
\item \concept{Setup Time ($t_{setup}$)}:
\begin{itemize}
\item Minimum time data must be stable before clock edge
\item Ensures proper data capture
\item Technology dependent
\end{itemize}
\vspace{0.3cm}
\item \concept{Hold Time ($t_{hold}$)}:
\begin{itemize}
\item Minimum time data must remain stable after clock edge
\item Prevents data corruption
\item Usually shorter than setup time
\end{itemize}
\vspace{0.3cm}
\item \concept{Propagation Delay ($t_{pd}$)}:
\begin{itemize}
\item Time from clock edge to output change
\item Determines maximum clock frequency
\item Includes internal delays
\end{itemize}
\vspace{0.3cm}
\item \concept{Clock-to-Q Delay ($t_{cq}$)}:
\begin{itemize}
\item Specific propagation delay measurement
\item Clock edge to output transition
\item Critical for timing analysis
\end{itemize}
\end{itemize}
\end{frame}

\begin{frame}{Timing Violations and Solutions}
\begin{itemize}
\item \concept{Setup Time Violation}:
\begin{itemize}
\item Data changes too close to clock edge
\item Results in metastability
\item Solutions: Slower clock, faster logic
\end{itemize}
\vspace{0.3cm}
\item \concept{Hold Time Violation}:
\begin{itemize}
\item Data changes too soon after clock edge
\item Causes incorrect data capture
\item Solutions: Add delay, buffer insertion
\end{itemize}
\vspace{0.3cm}
\item \concept{Clock Skew Effects}:
\begin{itemize}
\item Different clock arrival times
\item Can cause timing violations
\item Solutions: Clock tree synthesis, buffers
\end{itemize}
\vspace{0.3cm}
\item \concept{Metastability}:
\begin{itemize}
\item Unstable intermediate state
\item Occurs with timing violations
\item Solutions: Synchronizers, proper timing
\end{itemize}
\end{itemize}
\end{frame}

% Section 6: Memory Systems
\section{Random Access Memory (RAM)}

\begin{frame}{RAM Fundamentals}
\begin{columns}
\begin{column}{0.6\textwidth}
\concept{Definition}: Random Access Memory
\begin{itemize}
\item Any location accessible in constant time
\item Address-based data retrieval
\item Read and write operations
\item Parallel extension of 1-OS
\end{itemize}
\vspace{0.3cm}
\concept{Basic Structure}:
\begin{itemize}
\item $2^p$ memory locations
\item $p$-bit address input
\item Data input and output
\item Write enable control
\end{itemize}
\vspace{0.3cm}
\concept{Key Components}:
\begin{itemize}
\item Address decoder (DEMUX)
\item Memory cell array
\item Output multiplexer (MUX)
\item Control logic
\end{itemize}
\end{column}
\begin{column}{0.4\textwidth}
\begin{center}

\end{center}
\end{column}
\end{columns}
\end{frame}
\begin{frame}{RAM Operation and Timing}
\begin{columns}
\begin{column}{0.6\textwidth}
\concept{Read Operation}:
\begin{itemize}
\item Combinational process
\item Address selects memory cell
\item Data appears at output
\item Access time: $t_{ACC}$
\end{itemize}
\vspace{0.3cm}
\concept{Write Operation}:
\begin{itemize}
\item Requires write enable signal
\item Address must be stable
\item Data must meet timing requirements
\item More complex than read
\end{itemize}
\vspace{0.3cm}
\concept{Critical Timing Parameters}:
\begin{itemize}
\item $t_{ASU}$: Address setup time
\item $t_{DSU}$: Data setup time
\item $t_{DH}$: Data hold time
\item $t_{W}$: Write enable width
\item $t_{AH}$: Address hold time
\end{itemize}
\end{column}
\begin{column}{0.4\textwidth}
\begin{center}

\end{center}
\end{column}
\end{columns}
\end{frame}}
\begin{frame}{Multi-bit Word RAM}
\begin{columns}
\begin{column}{0.6\textwidth}
\concept{Word-Oriented Memory}:
\begin{itemize}
\item Store $m$-bit words instead of single bits
\item Parallel columns for each bit position
\item Shared address decoder
\item Independent data paths
\end{itemize}
\vspace{0.3cm}
\concept{Implementation}:
\begin{itemize}
\item One column per bit position
\item Each column contains $2^p$ memory cells
\item Common address decoding
\item Parallel read/write operations
\end{itemize}
\vspace{0.3cm}
\concept{Advantages}:
\begin{itemize}
\item Efficient data organization
\item Reduced addressing overhead
\item Better bandwidth utilization
\item Matches processor word size
\end{itemize}
\end{column}
\begin{column}{0.4\textwidth}
\begin{center}

\end{center}
\end{column}
\end{columns}
\end{frame}}

% Section 7: Synchronous Memory
\section{Synchronous Memory Systems}

\begin{frame}{Synchronous RAM (SRAM)}
\begin{columns}
\begin{column}{0.6\textwidth}
\concept{Motivation}: High-speed operation challenges
\begin{itemize}
\item GHz frequencies make timing difficult
\item Asynchronous timing constraints hard to meet
\item Need better controllability
\end{itemize>
\vspace{0.3cm}
\concept{Synchronous Solution}:
\begin{itemize}
\item All timing relative to clock edge
\item Setup and hold times defined for clock
\item More rigorous timing control
\item Better system integration
\end{itemize>
\vspace{0.3cm}
\concept{Benefits}:
\begin{itemize}
\item Simplified timing analysis
\item Better noise immunity
\item Easier system design
\item Predictable behavior
\end{itemize>
\end{column}
\begin{column}{0.4\textwidth}
\begin{center}

\end{center}
\end{column}
\end{columns}
\end{frame}

\begin{frame}{Pipelined Synchronous RAM}
\begin{itemize}
\item \concept{Performance Challenge}: Access time too high
\begin{itemize}
\item Memory access limits clock frequency
\item Need faster response time
\item Trade-off between speed and latency
\end{itemize}
\vspace{0.5cm}
\item \concept{Pipeline Solution}:
\begin{itemize}
\item Add output register
\item Fast access time (register delay)
\item One clock cycle latency
\item Higher system frequency possible
\end{itemize}
\vspace{0.3cm}
\item \concept{Design Considerations}:
\begin{itemize}
\item Latency vs. throughput trade-off
\item Pipeline depth optimization
\item Latency hiding techniques
\item System-level impact
\end{itemize}
\end{itemize}
\end{frame}

% Section 8: Register Files
\section{Register Files}

\begin{frame}{Register File Concept}
\begin{columns}
\begin{column}{0.6\textwidth}
\concept{Definition}: Battery of registers with multiple access ports
\begin{itemize}
\item Multiple simultaneous read ports
\item Single write port
\item Any register accessible each cycle
\item Critical processor component
\end{itemize}
\vspace{0.3cm}
\concept{Typical Configuration}:
\begin{itemize}
\item Two read ports (left and right)
\item One write port
\item $2^n$ registers of $m$ bits each
\item Independent address inputs
\end{itemize}
\vspace{0.3cm}
\concept{Applications}:
\begin{itemize}
\item CPU general-purpose registers
\item Pipeline stage buffers
\item Vector processor registers
\item DSP accumulator banks
\end{itemize}
\end{column}
\begin{column}{0.4\textwidth}
\begin{center}

\end{center}
\end{column}
\end{columns}
\end{frame}

\begin{frame}{Register File Implementation}
\begin{itemize}
\item \concept{Multi-Port Structure}:
\begin{itemize}
\item Two output multiplexers for reads
\item Address decoders for write selection
\item Write enable gating
\item Simultaneous operation capability
\end{itemize}
\vspace{0.5cm}
\item \concept{Interface Signals}:
\begin{itemize}
\item \texttt{leftAddr}: Left read port address
\item \texttt{rightAddr}: Right read port address
\item \texttt{destAddr}: Write destination address
\item \texttt{writeEnable}: Write control signal
\item \texttt{leftOut}, \texttt{rightOut}: Read data outputs
\item \texttt{in}: Write data input
\end{itemize}
\vspace{0.3cm}
\item \concept{Design Trade-offs}:
\begin{itemize}
\item Area vs. performance
\item Number of ports vs. complexity
\item Register count vs. access time
\end{itemize}
\end{itemize}
\end{frame}

% Section 9: Advanced Memory Concepts
\section{Advanced Memory Concepts}

\begin{frame}{Memory Hierarchy Considerations}
\begin{itemize}
\item \concept{Speed vs. Capacity Trade-off}:
\begin{itemize}
\item Faster memory is more expensive
\item Larger memory is slower
\item Need hierarchical approach
\end{itemize}
\vspace{0.5cm}
\item \concept{Memory Types by Speed}:
\begin{itemize}
\item \textbf{Registers}: Fastest, smallest capacity
\item \textbf{Cache}: Fast, medium capacity
\item \textbf{Main Memory}: Medium speed, large capacity
\item \textbf{Storage}: Slowest, largest capacity
\end{itemize}
\vspace{0.3cm}
\item \concept{Design Principles}:
\begin{itemize}
\item Locality of reference
\item Temporal and spatial locality
\item Caching strategies
\item Prefetching techniques
\end{itemize}
\end{itemize}
\end{frame}

\begin{frame}{Memory Interface Standards}
\begin{itemize}
\item \concept{SRAM Interfaces}:
\begin{itemize}
\item Synchronous operation
\item Clock-based timing
\item Burst mode capability
\item Pipeline-friendly
\end{itemize}
\vspace{0.3cm}
\item \concept{DRAM Interfaces}:
\begin{itemize}
\item Row/column addressing
\item Refresh requirements
\item Higher density
\item Complex timing
\end{itemize}
\vspace{0.3cm}
\item \concept{Emerging Technologies}:
\begin{itemize}
\item Non-volatile memories
\item 3D memory structures
\item Processing-in-memory
\item Quantum memory concepts
\end{itemize}
\end{itemize}
\end{frame}

% Section 10: Summary
\section{Chapter Summary}

\begin{frame}{Key Concepts Review}
\begin{itemize}
\item \concept{First-Order Systems (1-OS)}:
\begin{itemize}
\item Introduction of feedback loops
\item Memory capability through state storage
\item Foundation for sequential logic
\end{itemize}
\vspace{0.3cm}
\item \concept{Memory Elements}:
\begin{itemize}
\item Latches: Level-sensitive storage
\item Flip-flops: Edge-triggered storage
\item Timing constraints and analysis
\end{itemize}
\vspace{0.3cm}
\item \concept{Memory Systems}:
\begin{itemize}
\item RAM: Random access capability
\item Synchronous operation benefits
\item Register files for multi-port access
\end{itemize}
\end{itemize}
\end{frame}

\begin{frame}{Design Principles and Trade-offs}
\begin{itemize}
\item \concept{Timing Considerations}:
\begin{itemize}
\item Setup and hold time requirements
\item Clock skew and jitter effects
\item Synchronous vs. asynchronous design
\end{itemize}
\vspace{0.3cm}
\item \concept{Performance Optimization}:
\begin{itemize}
\item Pipeline techniques for speed
\item Latency vs. throughput trade-offs
\item Memory hierarchy design
\end{itemize}
\vspace{0.3cm}
\item \concept{System Integration}:
\begin{itemize}
\item Interface standardization
\item Timing closure methodology
\item Power and area considerations
\end{itemize}
\end{itemize}
\end{frame}

\begin{frame}{Connection to Higher-Order Systems}
\begin{itemize}
\item \concept{Building Blocks for 2-OS (Automata)}:
\begin{itemize}
\item State storage elements
\item Control signal generation
\item Sequence detection capability
\end{itemize}
\vspace{0.3cm}
\item \concept{Foundation for 3-OS (Processors)}:
\begin{itemize}
\item Register files for data storage
\item Pipeline registers for instruction flow
\item Memory systems for program/data storage
\end{itemize}
\vspace{0.3cm}
\item \concept{System-Level Impact}:
\begin{itemize}
\item Performance bottlenecks
\item Power consumption
\item Reliability and fault tolerance
\end{itemize}
\end{itemize}
\end{frame}

% Questions slide
\begin{frame}{Questions and Discussion}
\begin{center}
\Large
Questions?

\vspace{1cm}

\normalsize
\textbf{Discussion Topics:}
\begin{itemize}
\item How does feedback create memory in digital circuits?
\item What are the advantages of edge-triggered over level-sensitive storage?
\item Why is synchronous design preferred for high-speed systems?
\item How do timing constraints affect system performance?
\end{itemize}

\vspace{0.5cm}
\textbf{Next:} Chapter 2.10 - Automata: Second-Order Systems (2-OS)
\end{center}
\end{frame}

\end{document}

